% GENERATED FILE. Edit canonical lessons, then run npm run generate:books.
% canonical-source-hash: fnv1a64:94720205e123c676
% canonical-lessons: KA-C322-raja, KA-C322-rani, KA-C322-alu-rule, KA-C322-simhasana, KA-C322-kirita

\chapter{The King and the Queen}
\label{ch:ka-the-king-and-the-queen}
\hlchaptermodality{322}

\begin{chapteropening}
\textbf{By the end of this chapter:} \emph{I can say a king, a queen, to reign, a throne and a crown.}

Complete the last lesson of chapter 322: I can say a king, a queen, to reign, a throne and a crown.
\end{chapteropening}

\section[\texorpdfstring{\kn{ರಾಜ} (rāja)}{ರಾಜ (rāja)}]{\kn{ರಾಜ} (rāja) — a king}

\label{lesson:KA-C322-raja}



\begin{quote}
Before the new one: say the Kannada for a verdict, then the Kannada for a punishment.
\end{quote}

\subsection*{You'll want to know: \kn{ರಾಜ}}
\textbf{\kn{ರಾಜ}} — \emph{rāja} — \textquotedblleft{}a king\textquotedblright{}.

\begin{quote}
\kn{ನಾನು} \kn{ರಾಜ} \kn{ಆಗಿದ್ದರೆ} \kn{ಎಲ್ಲರಿಗೂ} \kn{ಮನೆ} \kn{ಕಟ್ಟಿಸುತ್ತಿದ್ದೆ}.
\end{quote}

\begin{quote}
\emph{nānu rāja āgiddare ellarigū mane kaṭṭisuttidde.}
\end{quote}

\begin{quote}
If I were king, I would have houses built for everyone.
\end{quote}

Plural: \kn{ರಾಜರು}, \emph{rājaru}.

The land a king rules is his \kn{ರಾಜ್ಯ}, and he lives in an \kn{ಅರಮನೆ}, a palace.

\begin{cousinweb}
\textbf{\kn{ರಾಜ}} (\emph{rāja}) is borrowed from Sanskrit \emph{rājan}, \textquotedblleft{}a chieftain, a king\textquotedblright{} (CDIAL 10679), in its compound form \emph{rāja}.
\end{cousinweb}

\begin{grammarlens}[title={What you've built}]
One more word for telling what happened.
\end{grammarlens}

\subsection*{Your turn}
\begin{itemize}
\raggedright
  \item \emph{Say it:} \emph{rāja}
  \item \emph{Say it:} \emph{rāja}, once more
  \item \emph{Say it:} \emph{nānu rāja āgiddare ellarigū mane kaṭṭisuttidde.} — If I were king, I would have houses built for everyone.
  \item \emph{Recall:} say the Kannada for a verdict, then the Kannada for a punishment, then say \emph{rāja} again
\end{itemize}

\begin{tcolorbox}[breakable,colback=teal!4,colframe=teal!35!black,title={Before you move on}]
What does \kn{ರಾಜ} mean? (\textquotedblleft{}A king\textquotedblright{}.) Say it once more.
\end{tcolorbox}
\section[\texorpdfstring{\kn{ರಾಣಿ} (rāṇi)}{ರಾಣಿ (rāṇi)}]{\kn{ರಾಣಿ} (rāṇi) — a queen}

\label{lesson:KA-C322-rani}



\begin{quote}
Before the new one: say the Kannada for a punishment, then the Kannada for a king.
\end{quote}

\subsection*{You'll want to know: \kn{ರಾಣಿ}}
\textbf{\kn{ರಾಣಿ}} — \emph{rāṇi} — \textquotedblleft{}a queen\textquotedblright{}.

\begin{quote}
\kn{ರಾಣಿ} \kn{ಅರಮನೆಯ} \kn{ತೋಟದಲ್ಲಿ} \kn{ನಡೆಯುತ್ತಿದ್ದಳು}.
\end{quote}

\begin{quote}
\emph{rāṇi aramaneya tōṭadalli naḍeyuttiddaḷu.}
\end{quote}

\begin{quote}
The queen was walking in the palace garden.
\end{quote}

Plural: \kn{ರಾಣಿಯರು}, \emph{rāṇiyaru}.

\begin{cousinweb}
\textbf{\kn{ರಾಣಿ}} (\emph{rāṇi}) goes back to Sanskrit \emph{rājñī}, \textquotedblleft{}a queen\textquotedblright{} (CDIAL 10692), the feminine of \emph{rājan}; the spoken languages of the north smoothed it to \emph{rāṇī}, and Kannada took that form.
\end{cousinweb}

\begin{grammarlens}[title={What you've built}]
One more word for telling what happened.
\end{grammarlens}

\subsection*{Your turn}
\begin{itemize}
\raggedright
  \item \emph{Say it:} \emph{rāṇi}
  \item \emph{Say it:} \emph{rāṇi}, once more
  \item \emph{Say it:} \emph{rāṇi aramaneya tōṭadalli naḍeyuttiddaḷu.} — The queen was walking in the palace garden.
  \item \emph{Recall:} say the Kannada for a punishment, then the Kannada for a king, then say \emph{rāṇi} again
\end{itemize}

\begin{tcolorbox}[breakable,colback=teal!4,colframe=teal!35!black,title={Before you move on}]
What does \kn{ರಾಣಿ} mean? (\textquotedblleft{}A queen\textquotedblright{}.) Say it once more.
\end{tcolorbox}
\section[\texorpdfstring{\kn{ಆಳು} (āḷu)}{ಆಳು (āḷu)}]{\kn{ಆಳು} (āḷu) — to reign; to rule}

\label{lesson:KA-C322-alu-rule}



\begin{quote}
Before the new one: say the Kannada for a king, then the Kannada for a queen.
\end{quote}

\subsection*{You'll want to know: \kn{ಆಳು}}
\textbf{\kn{ಆಳು}} — \emph{āḷu} — \textquotedblleft{}to reign; to rule\textquotedblright{}.

\begin{quote}
\kn{ಆ} \kn{ರಾಜ} \kn{ನಲವತ್ತು} \kn{ವರ್ಷ} \kn{ಈ} \kn{ನಾಡನ್ನು} \kn{ಆಳಿದನು}.
\end{quote}

\begin{quote}
\emph{ā rāja nalavattu varṣa ī nāḍannu āḷidanu.}
\end{quote}

\begin{quote}
That king ruled this land for forty years.
\end{quote}

Past (I …): \kn{ಆಳಿದೆ}, \emph{āḷide}.

What is ruled takes -\kn{ಅನ್ನು}. A different \kn{ಆಳು}, a noun, means a servant or a worker.

\begin{cousinweb}
\textbf{\kn{ಆಳು}} (\emph{āḷu}) is native Dravidian: DEDR 5157 lists Kannada \emph{āḷ}, \textquotedblleft{}to possess, govern, rule\textquotedblright{}, beside Tamil \emph{āḷ}, \textquotedblleft{}to rule, reign over\textquotedblright{}.
\end{cousinweb}

\begin{grammarlens}[title={What you've built}]
One more word for telling what happened.
\end{grammarlens}

\subsection*{Your turn}
\begin{itemize}
\raggedright
  \item \emph{Say it:} \emph{āḷu}
  \item \emph{Say it:} \emph{āḷu}, once more
  \item \emph{Say it:} \emph{ā rāja nalavattu varṣa ī nāḍannu āḷidanu.} — That king ruled this land for forty years.
  \item \emph{Recall:} say the Kannada for a king, then the Kannada for a queen, then say \emph{āḷu} again
\end{itemize}

\begin{tcolorbox}[breakable,colback=teal!4,colframe=teal!35!black,title={Before you move on}]
What does \kn{ಆಳು} mean? (\textquotedblleft{}To reign; to rule\textquotedblright{}.) Say it once more.
\end{tcolorbox}
\section[\texorpdfstring{\kn{ಸಿಂಹಾಸನ} (siṃhāsana)}{ಸಿಂಹಾಸನ (siṃhāsana)}]{\kn{ಸಿಂಹಾಸನ} (siṃhāsana) — a throne}

\label{lesson:KA-C322-simhasana}



\begin{quote}
Before the new one: say the Kannada for a queen, then the Kannada for to reign.
\end{quote}

\subsection*{You'll want to know: \kn{ಸಿಂಹಾಸನ}}
\textbf{\kn{ಸಿಂಹಾಸನ}} — \emph{siṃhāsana} — \textquotedblleft{}a throne\textquotedblright{}.

\begin{quote}
\kn{ರಾಜ} \kn{ಸಿಂಹಾಸನದ} \kn{ಮೇಲೆ} \kn{ಕುಳಿತಿದ್ದನು}.
\end{quote}

\begin{quote}
\emph{rāja siṃhāsanada mēle kuḷitiddanu.}
\end{quote}

\begin{quote}
The king was sitting on the throne.
\end{quote}

The lion seat, as \kn{ಸಿಂಹ} is a lion. \kn{ಸಿಂಹಾಸನ} \kn{ಏರು}, to climb the throne, is to become king.

\begin{cousinweb}
\textbf{\kn{ಸಿಂಹಾಸನ}} (\emph{siṃhāsana}) is borrowed from Sanskrit \emph{siṃhāsana}, \textquotedblleft{}the lion's seat, the king's seat, a throne\textquotedblright{} (Monier-Williams), from \emph{siṃha} \textquotedblleft{}a lion\textquotedblright{} and \emph{āsana} \textquotedblleft{}a seat\textquotedblright{}.
\end{cousinweb}

\begin{grammarlens}[title={What you've built}]
One more word for telling what happened.
\end{grammarlens}

\subsection*{Your turn}
\begin{itemize}
\raggedright
  \item \emph{Say it:} \emph{siṃhāsana}
  \item \emph{Say it:} \emph{siṃhāsana}, once more
  \item \emph{Say it:} \emph{rāja siṃhāsanada mēle kuḷitiddanu.} — The king was sitting on the throne.
  \item \emph{Recall:} say the Kannada for a queen, then the Kannada for to reign, then say \emph{siṃhāsana} again
\end{itemize}

\begin{tcolorbox}[breakable,colback=teal!4,colframe=teal!35!black,title={Before you move on}]
What does \kn{ಸಿಂಹಾಸನ} mean? (\textquotedblleft{}A throne\textquotedblright{}.) Say it once more.
\end{tcolorbox}
\section[\texorpdfstring{\kn{ಕಿರೀಟ} (kirīṭa)}{ಕಿರೀಟ (kirīṭa)}]{\kn{ಕಿರೀಟ} (kirīṭa) — a crown}

\label{lesson:KA-C322-kirita}



\begin{quote}
Before the new one: say the Kannada for to reign, then the Kannada for a throne.
\end{quote}

\subsection*{You'll want to know: \kn{ಕಿರೀಟ}}
\textbf{\kn{ಕಿರೀಟ}} — \emph{kirīṭa} — \textquotedblleft{}a crown\textquotedblright{}.

\begin{quote}
\kn{ರಾಣಿಯ} \kn{ಕಿರೀಟ} \kn{ಚಿನ್ನದ್ದು}.
\end{quote}

\begin{quote}
\emph{rāṇiya kirīṭa cinnaddu.}
\end{quote}

\begin{quote}
The queen's crown is made of gold.
\end{quote}

\begin{cousinweb}
\textbf{\kn{ಕಿರೀಟ}} (\emph{kirīṭa}) is borrowed from Sanskrit \emph{kirīṭa}, \textquotedblleft{}a diadem, a crest\textquotedblright{} (CDIAL 3176).
\end{cousinweb}

\begin{grammarlens}[title={What you've built}]
One more word for telling what happened.
\end{grammarlens}

\subsection*{Your turn}
\begin{itemize}
\raggedright
  \item \emph{Say it:} \emph{kirīṭa}
  \item \emph{Say it:} \emph{kirīṭa}, once more
  \item \emph{Say it:} \emph{rāṇiya kirīṭa cinnaddu.} — The queen's crown is made of gold.
  \item \emph{Recall:} say the Kannada for to reign, then the Kannada for a throne, then say \emph{kirīṭa} again
\end{itemize}

\begin{tcolorbox}[breakable,colback=teal!4,colframe=teal!35!black,title={Before you move on}]
What does \kn{ಕಿರೀಟ} mean? (\textquotedblleft{}A crown\textquotedblright{}.) Say it once more.
\end{tcolorbox}
